\tikzset{
    ddfv/.style={
        scale=1.5,
        every path/.style={thick,line join=round, cap=round},
        ptsq/.style 2 args={
          rectangle,
          fill=##1,
          inner sep=2pt,
          label={[text=##1]##2}
        },
        ptsqbord/.style 2 args={
          rectangle,
          draw=##1,
          inner sep=2pt,
          label={[text=##1]##2}
        },
        ptcl/.style 2 args={
          circle,
          fill=##1,
          inner sep=1.5pt,
          label={[text=##1]##2}
        },
        ptclbord/.style 2 args={
          circle,
          draw=##1,
          inner sep=1.5pt,
          label={[text=##1]##2}
        },
        mailleprimale/.style={
            myblue,
            fill=mylightblue
        },
        mailledualeint/.style={
            myred,
            fill=mylightred
        },
        mailledualebord/.style={
            myred,
            fill=mylightred!50!white
        },
        maillediamantint/.style={
            mygreen,
            fill=mylightgreen
        },
        maillediamantbord/.style={
            mygreen,
            fill=mylightgreen!50!white
        }
    }
}

\begin{tikzpicture}[
  ddfv,
  face/.style={
    draw=myblue,
    line join=round
  },
]

% Longueur d'une arête
\def\arete{16mm}

% ------------------------------------------------------------
% \Triangle{numéro}{contenu côté 1}{contenu côté 2}{contenu côté 3}
%
% Dans le repère local, les sommets sont :
%   A_i = (0,0),
%   B_i = (\arete,0),
%   C_i = (60:\arete).
%
% Le centre de gravité est enregistré sous le nom G<i>.
% ------------------------------------------------------------
\newcommand{\Triangle}[4]{%
  % Sommets de la face
  \coordinate (A#1) at (0,0);
  \coordinate (B#1) at (\arete,0);
  \coordinate (C#1) at (60:\arete);

  % Centre de gravité
  \coordinate (G#1) at
    (barycentric cs:A#1=1,B#1=1,C#1=1);

  \coordinate (M#1) at ($(A#1)!0.5!(B#1)$);
  \coordinate (MB#1) at ($(A#1)!0.5!(C#1)$);
  \coordinate (MT#1) at ($(B#1)!0.5!(C#1)$);

  % Premier côté : C_i -- B_i
  \begin{scope}[
    shift={(60:\arete)},
    rotate=-60
  ]
    #2
  \end{scope}

  % Deuxième côté : B_i -- A_i
  \begin{scope}[
    xshift=\arete,
    rotate=180
  ]
    #3
  \end{scope}

  % Troisième côté : A_i -- C_i
  \begin{scope}[
    rotate=60
  ]
    #4
  \end{scope}
}

% Une arête libre ne porte aucun triangle supplémentaire.
\newcommand{\Libre}{}

% ------------------------------------------------------------
% Patron de l'icosaèdre.
%
% La structure d'emboîtement est celle de
% \tikzfoldingicosahedron dans tikzlibraryfolding.code.tex.
% ------------------------------------------------------------

\Triangle{1}
{
    \Triangle{2}
    {
        \Triangle{3}
        {\Libre}
        {\Libre}
        {\Libre}
    }
    {\Libre}
    {
        \Triangle{4}
        {
            \Triangle{5}
            {
                \Triangle{6}
                {\Libre}
                {\Libre}
                {
                    \Triangle{7}
                    {\Libre}
                    {\Libre}
                    {
                        \Triangle{8}
                        {
                            \Triangle{9}
                            {
                                \Triangle{11}
                                {
                                    \Triangle{12}
                                    {\Libre}
                                    {\Libre}
                                    {
                                        \Triangle{13}
                                        {\Libre}
                                        {\Libre}
                                        {
                                            \Triangle{14}
                                            {\Libre}
                                            {\Libre}
                                            {
                                                \Triangle{15}
                                                {\Libre}
                                                {\Libre}
                                                {\Libre}
                                            }
                                        }
                                    }
                                }
                                {\Libre}
                                {\Libre}
                            }
                            {\Libre}
                            {\Libre}  
                        }
                        {\Libre}
                        {
                            \Triangle{10}
                            {
                                \Triangle{16}
                                {
                                    \Triangle{17}
                                    {\Libre}
                                    {\Libre}
                                    {
                                        \Triangle{18}
                                        {\Libre}
                                        {\Libre}
                                        {
                                            \Triangle{19}
                                            {\Libre}
                                            {\Libre}
                                            {
                                                \Triangle{20}
                                                {\Libre}
                                                {\Libre}
                                                {\Libre}
                                            }
                                        }
                                    }
                                }
                                {\Libre}
                                {\Libre}
                            }
                            {\Libre}
                            {\Libre}
                        }
                    }
                }
            }
            {\Libre}
            {\Libre}
        }
        {\Libre}
        {\Libre}
    }
    }
{\Libre}
{\Libre}

\fill[mylightred] (G3) -- (G2) -- (G4) -- (G5) -- (G6) -- (MT6) -- (B6) -- (MB3) -- cycle;

\foreach \i in {1,2,...,20}{
    \draw[myred] (G\i) -- (M\i);
    \draw[myred] (G\i) -- (MB\i);
    \draw[myred] (G\i) -- (MT\i);

    \draw[black,semithick] (G\i) -- (A\i);
    \draw[black,semithick] (G\i) -- (B\i);
    \draw[black,semithick] (G\i) -- (C\i);

    % Triangle
    \draw[face]
      (A\i) -- (B\i) -- (C\i) -- cycle;

    \node[ptsq={myblue}{}] at (A\i) {};
    \node[ptsq={myblue}{}] at (B\i) {};
    \node[ptsq={myblue}{}] at (C\i) {};

    \node[ptcl={myred}{}] at (G\i) {};

    \node[ptcl={black}{}] at (M\i) {};
    \node[ptcl={black}{}] at (MB\i) {};
    \node[ptcl={black}{}] at (MT\i) {};
}

\node[myblue, below right] at (B6) {\contour{white}{$x_{\mailleduale{K}}$}};
\node[below] at (MT6) {$m_\sigma$};
\node[myred, above=3pt] at (G6) {\contour{white}{$x_{\primalplat{K}}$}};

% \node[myblue, below] at (B9) {$x_{\mailleduale{K}}$};

\end{tikzpicture}